\cotefile{157-1}
\scanpage{182}\mspage{38}
% archive sheet number: 182

Ainsi, les foncteurs familiers, \deleted{pour} une \deleted{flèche}

\[
u:I'\unsure{\longrightarrow|\mathbin{-}}I
\]

dans $\mathrm{Cat}$

\diagram{$(\unsure{3.7.})\quad\left\{\vphantom{\rule{0pt}{4em}}\right.$\begin{tikzcd}[ampersand replacement=\&,column sep=large]
I'{}^{\wedge} \arrow[r,shift left=4ex,"{u_{!}}"] \arrow[r,shift right=4ex,"{u_*}"] \& I^{\wedge} \arrow[l,"{u^*}"']
\end{tikzcd}$,$}

s'interprètent comme les foncteurs de \unsure{même nom|mêmes noms} pour le dérivateur $\mathrm{HOT}_{W_0}$, ou \inserted{plus} directement, pour le dérivateur (qui pouvait \unsure{justement|faussement|jusqu'ici} \unsure{être} considéré comme \unsure{canularesque}) $\mathbb{D}_{(\mathrm{Ens})}$.

Cela donne un sens à \deleted{\unsure{une autre}} la cohorte des notions cohomologiques associées à un dérivateur, ou plus particulièrement à un localiseur fondamental $W$ sur $\mathrm{Cat}$ (donnant naissance au dérivateur $\mathrm{HOT}_W$) \unsure{;} \unsure{On|les} \deleted{flèches} \deleted{$W_0$-équivalences \unsure{:}} les \unsure{désignera|désigner} par les noms de $W_0$-\unsure{foncté} (au lieu de $\mathrm{HOT}_{W_0}$-foncté, ou $\mathbb{D}_{\mathrm{Ens}}$-foncté), \deleted{\unsure{$\mathbb{D}$}} \deleted{\unsure{tels}} notamment pour une flèche $u:I'\longrightarrow I$ de $\mathrm{Cat}$ telle la notion

\[
\left\{\begin{aligned}
& W_0\text{-équivalence}\quad(\Longleftrightarrow\pi_0(u)\text{ bijectif})\\
& \vert\ \ W_0\text{-asphérique}\quad(\Longleftrightarrow\text{ les }I'_{\deleted{\unsure{u}}/i}\text{ pour }i\in\mathrm{Ob}\,I\text{ sont}\\[-2pt]
&\hspace{13em}0\text{-connexes}\\
& \vert\ \ W_0\text{-coasphérique}\quad(\Longleftrightarrow\text{ les }{}_{i\backslash}I'\text{ pour }i\in\mathrm{Ob}\,I\text{ sont}\\[-2pt]
&\hspace{13em}0\text{-connexes})\\
& W_0\text{-propre}\quad(\Longleftrightarrow\text{les inclusions }I'_i\longrightarrow I'_{/i}\ (\unsure{i'|i}\in\mathrm{Ob}\,\unsure{I|I'})\\
&\hspace{7em}\text{sont }W_0\text{-coasphériques, i.e. les fibres}\\
&\hspace{7em}\text{des }\deleted{\unsure{l}}\text{colocalisés }{}_{i'\backslash}I'_{\deleted{\unsure{/i}}}\longrightarrow{}_{i\backslash}I_{\deleted{\unsure{/i}}}\ (\unsure{i'=ui|i=ui'})\\
&\hspace{7em}\text{sont }0\text{-connexes}
\end{aligned}\right. \tag*{(3.8)}
\]

\margin{Ex : les Cat-\deleted{\unsure{fibrations}} cofibrations,}
